\documentclass[twocolumn, a4paper]{UECIEresume}
\usepackage[dvipdfmx]{graphicx}
\usepackage{amsmath}
\usepackage{txfonts}
\usepackage{url}

\title{Cross-Lingual Zero-Shot TTSにおける話者性を保持した日本語音声合成の検討}
\date{2026 年 09 月 16 日}
\affiliation{I類 メディア情報学 プログラム}
\supervisor{髙木一幸 助教}
\studentid{2311093}
\author{駒月柊平}
\headtitle{2026 年度 I類 卒研中間発表}

\begin{document}
\maketitle

\section{はじめに}

近年，短いプロンプト音声を用いて，学習時に含まれない話者の音声を
追加学習なしで合成するzero-shot text-to-speech（TTS）が研究されている
\cite{Wang2023VALLE}．
さらに，プロンプト音声とは異なる言語でも話者性を保持した音声を生成する
cross-lingual zero-shot TTS（以下，cross-lingual TTS）が提案されている
\cite{Zhang2023VALLEX}．
これにより，対象話者が実際には話していない言語でも，
その話者らしい音声を合成することが期待される．

一方，cross-lingual TTSでは，対象言語として正確で自然な発話を生成することと，
プロンプト音声の話者性を保持することの両方が求められる．近年のcross-lingual TTSでは，
言語方向によって発話内容の正確性や話者類似性が異なることが報告されている
\cite{Xu2026XVoice}．

本研究に先立ち，同一の推論条件で英語および日本語を対象言語とした
予備的な生成を比較したところ，日本語を対象言語とする条件では，
英語を対象言語とする条件に比べて，自然性および話者類似性が低いと
感じられる例が観察された．特に，プロンプト音声と合成音声を
同一話者の発話として知覚しにくい例がみられた．

これを踏まえ，対象言語を日本語とし，cross-lingual TTSにおける
発話内容の正確性，自然性，および話者類似性の品質低下の有無およびその要因を，
客観評価および主観評価の両面から検討する．その結果に基づき，
話者性を保持した日本語音声合成の改善を目指す．
まず既存モデルの生成過程を構成要素ごとに分析し，条件間の品質差が生じる要因を
調査する．その結果に基づいて改善対象と手法を決定し，既存モデルとの比較により
有効性を検証する．

\section{研究方針}

本節では，cross-lingual条件で生じる話者性の変化の要因を，
生成過程の構成要素ごとに分析する方針を述べる．
主な分析対象にはCosyVoice 3\cite{Du2025CosyVoice3}を用いる．CosyVoice 3は，
テキストから離散音声トークンを予測するLLMベースのtext-to-token部と，
予測されたトークンおよびプロンプト音声から抽出した話者埋め込みを条件として
音響特徴を生成する条件付きflow matchingベースのtoken-to-speech部から構成される．

話者性の変化がtext-to-token部，token-to-speech部，または両者のいずれに
起因するかを明らかにする．cross-lingual条件では，プロンプト音声から抽出した
speech tokenはLLMの入力に用いられず，話者埋め込みのみがtoken-to-speech部の
条件として用いられる．そこで，LLMがspeech tokenを予測する通常生成と，
生成対象と同一話者・同一発話の実音声から抽出したspeech tokenを用いる条件を
比較する．これにより，text-to-token部による影響とtoken-to-speech部による影響を
分けて検討する．さらに，話者埋め込みおよびspeech tokenを個別に変更する条件を設け，
話者性の変化に寄与する要因を分析する．

これらの分析は，改善対象となるモジュールを選定するために行うものであり，
特定の改善手法の有効性を直接示すものではない．改善対象を特定した後は，
対象モジュールには，LoRA（Low-Rank Adaptation）やadapterなどの
パラメータ効率的な追加学習手法を適用する．その上で，学習に用いる
データセットを決定する．

\section{データセット}

予備分析には，日本語と英語の音声を同一条件で比較するため，1名のバイリンガル
話者による日英パラレル音声を収録したJECS（Japanese-English Code-Switching
Speech Corpus）\cite{JECS}を用いる．
ただし，収録話者が1名であるため，得られた傾向の話者一般への一般化には限界がある．

複数話者における分析には，100名の日本語話者を収録したJVS（Japanese Versatile
Speech）\cite{Takamichi2019JVS}と，109名の英語話者を収録したVCTK
\cite{VeauxVCTK}を用いる．
% これらはspeech tokenと話者条件の関係を各言語内で
% 調べるために用いるが，同一人物の日英音声ではないため，コーパス間のSIMの
% 絶対値を言語差として比較しない．

提案手法の学習に用いるデータは，分析により更新対象を決定した後，話者数，
収録時間，録音品質，利用条件を考慮して選定する．学習を行う場合は，学習用，
検証用，評価用の話者を分離し，zero-shot評価の話者を学習データに含めない．

% 「最後の改善対象を特定した後は」の一文が長い、冗長
% 「学習に用いるデータセットを決定する」という表現がつながっていない
% 
% 

\section{評価方法}

CosyVoice 3をはじめとする既存のcross-lingual TTSモデルと，
本研究で検討する改善手法を，発話内容の正確性，自然性，
話者類似性の3観点から比較評価する．
% 条件間では同一のプロンプト音声と入力文を用い，
% 自己回帰生成については複数の乱数seedによる平均とばらつきを報告する．

\subsection{客観評価}

発話内容の正確性は，合成音声を多言語音声認識モデルWhisper large-v3
\cite{Radford2023Whisper,OpenAI2023WhisperLargeV3}で文字起こしして評価する．英語には単語誤り率
（Word Error Rate: WER），日本語には文字誤り率（Character Error Rate: CER）を
用いる．置換，削除，挿入の数をそれぞれ$N_{\mathrm{S}}$，$N_{\mathrm{D}}$，
$N_{\mathrm{I}}$，正解文の単語数または文字数を$N$として，誤り率を
\begin{equation}
  \mathrm{Error\ Rate}
  = \frac{N_{\mathrm{S}}+N_{\mathrm{D}}+N_{\mathrm{I}}}{N}
\end{equation}
と定義する．
% 日本語では，漢字と仮名などの表記差を発音誤りとして数える影響を
% 確認するため，正解文と認識結果を読み表記に統一したKana-CERも補助的に報告する．
% 全条件で大文字・小文字，空白，句読点，数字などの正規化規則を統一する．また，
% 実音声にも同じ認識処理を適用し，ASR自体の誤りを基準値として示す．

話者類似性にはSpeaker Similarity（SIM）を用いる．生成モデルの条件付けには
使用されていない外部の話者照合モデルWeSpeaker
\cite{Wang2023WeSpeaker}により，プロンプト音声と合成音声から話者埋め込み
$\boldsymbol{e}_{\mathrm{p}}$と$\boldsymbol{e}_{\mathrm{s}}$を抽出し，
\begin{equation}
  \mathrm{SIM}
  = \frac{\boldsymbol{e}_{\mathrm{p}}^{\mathsf{T}}
  \boldsymbol{e}_{\mathrm{s}}}
  {\lVert\boldsymbol{e}_{\mathrm{p}}\rVert_2
  \lVert\boldsymbol{e}_{\mathrm{s}}\rVert_2}
\end{equation}
によりコサイン類似度を求める．SIMは，当該埋め込み空間における類似性の
代理指標である．話者照合モデルの言語依存性を確認するため，同一話者の
実音声間についても同一言語・異言語のSIMを算出する．自然性の自動評価には，
音声から知覚品質に関するMean Opinion Score（MOS）を推定する
UTMOS\cite{Saeki2022UTMOS}を用いる．
% ただし，MOS推定モデルは，学習時と異なる聴取実験の条件や音声合成方式に対して
% 予測精度が低下し得る\cite{Huang2022VoiceMOS}．したがって，UTMOSは条件間の
% 傾向を把握するための補助指標として用い，自然性は聴取実験によっても評価する．

\subsection{主観評価}

客観指標のみでは知覚的な自然性や話者らしさを十分に評価できないため，
聴取実験を行う．
% ITU-T P.800\cite{ITUP800}を参考に，
評価者数は10名以上を目安として確保する．
自然性はMOSを用い，「非常に不自然」から「非常に自然」までの5段階で
評価する．話者類似性はSimilarity Mean Opinion Score（SMOS）を用い，
参照音声と合成音声を提示して，「全く似ていない」から「非常に似ている」までの
5段階で評価する．日本語音声については，日本語母語話者を評価者とする．

改善手法と基盤モデルの差を直接比較する必要がある場合は，提示順を
無作為化した比較評価を行う．自然性など評価軸を一つずつ提示し，基盤モデルに
対する改善手法の相対評価を$-3$から$+3$の7段階で求める．
% 評価音声数，評価者数，
% 提示順序，評価尺度のラベルを固定し，平均値と信頼区間を報告する．

\section{おわりに}

本研究では，cross-lingual TTSにおける話者性を保持した日本語音声合成を目的とし，
既存モデルの生成過程を構成要素ごとに分析する．現在は，プロンプト言語と
生成言語を分けた通常生成，実音声のspeech tokenを用いた生成，および話者条件の
変更実験により，改善対象を選定するための検討を進めている．

今後は，複数話者を対象として分析結果の再現性を確認し，
その結果に基づいて改善対象，改善手法，および学習データを決定する．
その後，提案手法について，発話内容の正確性，自然性，および話者類似性を
客観評価と主観評価の両面から検証する．

{\small
\bibliographystyle{junsrt}
\bibliography{template}
}
\end{document}
